\section{Заключение}
\label{sec:conclusion}

Настоящата дипломна работа представи проектирането, реализацията и анализа на \ORMName{} библиотека за C++, ориентирана към работа както с релационни, така и с нерелационни бази данни. В центъра на разработката стои идеята, че query структурата, логическият модел на данните и съществена част от ограниченията върху употребата на backend-ите могат да бъдат описани и валидирани чрез типовата система на езика.

С това работата адресира един конкретен проблем от съвременната практика: разминаването между силно типизиран приложен код и хетерогенните модели за съхранение, които този код трябва да използва. Класическите ORM решения са най-силни в рамките на релационния свят, а чисто driver-oriented подходите често прехвърлят твърде голяма част от коректността към runtime. Настоящата разработка показа, че между тези два полюса съществува жизнеспособна архитектурна позиция --- compile-time ORM слой, който е едновременно практичен, формално структуриран и достатъчно честен по отношение на backend различията.

\subsection{Обобщение на решената задача}

Поставената задача не беше просто да се изгради още един удобен query builder. Целта беше да се проектира библиотека, в която:

\begin{enumerate}[label=\arabic*.]
    \item логическият модел на данните се описва чрез C++ типове;
    \item заявките се представят като типизиран IR, а не като низове;
    \item backend-ите се интегрират чрез формален connector contract;
    \item различията между storage моделите се управляват чрез capability и constraint механизми;
    \item коректността се подкрепя едновременно от compile-time проверки и от тестова инфраструктура.
\end{enumerate}

В рамките на дипломната работа всяка от тези подцели беше адресирана с конкретни архитектурни и implementation решения. Entity слоят въведе типизирани свойства и отношения, политика за съхранение чрез \code{store\_as} и traits за именуване на логическите контейнери. Query builder-ът беше организиран около типизиран intermediate representation. Connector слойът беше формализиран чрез \code{connector\_trait<DB>}, а verification стратегията бе подкрепена от unit и integration тестове за множество backend-и.

\begin{figure}[H]
\centering
\begin{tikzpicture}[node distance=1.35cm and 1.15cm]
    \node[ormaccent] (model) {Логически model layer\\entity types + relationships};
    \node[ormblock, right=of model] (ir) {Typed query IR\\predicates, projections, placeholders};
    \node[ormblock, below=of $(model)!0.5!(ir)$] (connector) {Connector contract\\capabilities + materialization};
    \node[ormblock, below left=of connector] (tests) {Verification\\unit + integration tests};
    \node[ormblock, below right=of connector] (backends) {Multi-backend\\execution};

    \draw[ormline] (model) -- (ir);
    \draw[ormline] (model) -- (connector);
    \draw[ormline] (ir) -- (connector);
    \draw[ormline] (connector) -- (tests);
    \draw[ormline] (connector) -- (backends);
\end{tikzpicture}
\caption{Синтез на основната конструкция на разработената библиотека}
\label{fig:thesis_synthesis}
\end{figure}

Фигура~\ref{fig:thesis_synthesis} обобщава цялостната логика на дипломната работа. Тя показва, че библиотеката не е съвкупност от отделни техники, а последователна система, в която model layer, query IR, connector contract, verification и multi-backend execution се поддържат взаимно.

\subsection{Обобщение на постигнатите резултати}

Разработката показа, че C++ предоставя достатъчно силни механизми --- шаблони, \code{constexpr}, \code{concepts} и trait специализации --- за изграждане на ORM слой, в който заявките не се описват като низове. Вместо това те се представят чрез типизиран query IR, който може да бъде анализиран на compile-time и материализиран по различен начин според backend-а.

На равнище модел на данните библиотеката демонстрира как \code{property}, \code{relationship}, \code{store\_as} и \code{table\_name\_trait} могат да играят ролята на единен логически слой. Този слой не е ограничен до релационната интерпретация, а служи като общо описание на данните и връзките между тях. Именно това позволява една и съща application-level структура да бъде пренасяна към таблици, документи, ключове и стойности, графови възли или wide-column редове.

На равнище архитектура connector trait-ът е формализиран като централен механизъм за разширяемост. Именно чрез него библиотеката адаптира един и същ query IR към SQL, BSON-подобни филтри, Redis команди, Cypher и CQL. Capability механизмът допълва тази архитектура, като прави недопустимите операции видими още при компилация.

На равнище верификация разработката се опира на систематичен набор от unit и integration тестове, които покриват entity слоя, query builder-а, prepared query механизма, миграциите, нишковата безопасност и конкретните конектори. По този начин архитектурните твърдения са подкрепени с реални доказателства от кода и тестовата инфраструктура.

\subsection{Научни и приложни приноси}

Първият принос на работата е демонстрацията, че practically useful \ORMName{} библиотека може да бъде реализирана в стандартен C++. Това се постига без външна кодогенерация, без виртуален connector интерфейс и без принуждаване на потребителя да напуска идиоматичния C++ модел.

Вторият принос е в изясняването на връзката между C++ entity описанието и физическата организация на данните. Показано е, че C++ моделът може да бъде третиран като първичен носител на логическата схема, докато конкретният backend определя начина на материализация. Тази постановка има стойност не само за ORM библиотеките, а и по-общо за системите, които се опитват да комбинират силна статична типизация с хетерогенни източници на данни.

Третият принос е формализирането на connector contract, който обединява разширяемост и строга compile-time дисциплина. Този contract прави възможно надграждането на библиотеката с нови backend-и, без да се разрушават вече съществуващите гаранции в query API-то. Точно тук работата се различава от подходи, които предлагат общ интерфейс, но скриват критичните разлики между backend-ите.

Четвъртият принос е показването, че verification-ът на подобна библиотека трябва да бъде многостепенен. Не е достатъчно само ``да има тестове'' или само ``да има compile-time checks''. Необходима е комбинация от двете: concept-level constraints, capability assertions, query rendering tests, live integration сценарии и migration/thread-safety проверки.

\begin{table}[H]
\centering
\small
\caption{Връзка между целите на дипломната работа и реализираните резултати}
\label{tab:goal_result_mapping}
\begin{tabularx}{\textwidth}{L{3.1cm}L{3.2cm}Y}
\toprule
\textbf{Цел} & \textbf{Реализиран механизъм} & \textbf{Доказателствена опора} \\
\midrule
Типизиран модел на данните & \code{property}, \code{relationship}, \code{table\_name\_trait} & entity headers и unit тестовете за property/relationship \\
Типизиран query слой & query builder-и и typed IR & query headers, \path{tests/unit/test_query.cpp}, \path{tests/integration/test_mockdb.cpp} \\
Разширяем multi-backend слой & connector contract и capability tags & connector headers и backend-specific unit/live тестове \\
Schema-aware посока & \code{migrate<DB>} и DDL hooks & \path{tests/unit/test_schema_migration.cpp} \\
Безопасна operational употреба & thread-safety и transaction helper-и & \path{tests/unit/test_thread_safety.cpp} \\
\bottomrule
\end{tabularx}
\end{table}

Таблица~\ref{tab:goal_result_mapping} показва, че дипломната работа не остава на равнището на декларативен design document. За всяка основна цел има конкретен реализиран механизъм и поне един пряк технически корелат в хранилището.

\subsection{Основни ограничения и граници на обобщението}

Работата също така ясно показва, че multi-backend ORM абстракцията има граници. Не всеки модел за съхранение предлага едни и същи операции, а следователно общият query IR не може да бъде интерпретиран еднакво навсякъде. Именно затова capability и constraint механизмите са толкова важни. Те не отслабват библиотеката, а я правят коректна спрямо различните backend-и.

Допълнително, макар архитектурната основа за миграции, нишкова безопасност и по-ниско ниво на изпълнение вече да съществува, тези подсистеми остават естествено поле за по-нататъшно развитие и емпирично разширяване. Това означава, че работата трябва да бъде оценявана като завършена архитектурна и implementation демонстрация, но не и като последна дума по темата за compile-time достъп до данни.

От академична гледна точка е важно да се подчертае и още една граница. Настоящата разработка не доказва, че compile-time ORM е универсално най-добрият подход за всички приложения. Тя доказва нещо по-точно: че за определен клас системи --- системи с нужда от ясна статична структура, висока type safety и controlled multi-backend extensibility --- този подход е жизнеспособен и инженерно защитим.

\subsection{Значение за практиката и за бъдещите изследвания}

Практическата стойност на разработката се състои в това, че библиотеката вече осигурява стабилна основа за по-нататъшно инженерно развитие. Entity описанието, query IR-ът, connector contract-ът и verification стратегията са достатъчно ясно формулирани, за да позволят дисциплинирано добавяне на нови backend-и и operational възможности.

Изследователската стойност е свързана с по-общия въпрос за ролята на типовата система в моделирането на достъпа до данни. Работата показва, че type system-ът на C++ може да служи не само за описание на данни и алгоритми, но и за формализиране на архитектурни граници, допустимост на операции и преход между логически и физически слоеве. Това поставя библиотеката в по-широк контекст, който надхвърля ORM темата в тесния смисъл.

\subsection{Заключителни бележки}

Направеният анализ потвърждава, че типовата система на C++ е достатъчно мощна, за да служи не само за описване на данни и алгоритми, но и за формализиране на достъпа до персистентни системи. Комбинацията от статичен entity модел, query IR, trait-базиран connector слой и систематична верификация показва устойчив път към библиотеки, които са едновременно изразителни, безопасни и надграждаеми.

По този начин дипломната работа постига както инженерна, така и изследователска цел: изгражда работеща библиотека и едновременно формулира принципи, които могат да бъдат използвани при бъдещи разработки на compile-time слоеве за достъп до данни. Най-същественият извод е, че compile-time ORM подходът не е просто теоретично упражнение, а реална инженерна алтернатива за системи, които изискват строгост, яснота и разширяемост при работа с множество модели за съхранение.
